\thispagestyle{fancy}
\begin{center}
    \textbf{\Large TÓM TẮT LUẬN VĂN}
\end{center}

\vspace{0.5cm}

\noindent\textbf{Tên đề tài:} Hệ thống LMS tự động tạo Micro-Content \& Quiz
bằng Generative AI\\[4pt]
\noindent\textbf{Giảng viên hướng dẫn:} [Tên Giảng Viên]\\[4pt]
\noindent\textbf{Sinh viên thực hiện:} Phùng Xương Cận (2210348)

\vspace{0.7cm}

Việc chuyển đổi giáo trình, slide và tài liệu học thuật thành nội dung học tập
ngắn gọn, câu hỏi kiểm tra và tài nguyên ôn tập thường tiêu tốn nhiều thời
gian của giảng viên. Các công cụ AI phổ biến có khả năng tóm tắt hoặc sinh câu
hỏi, nhưng thường thiếu quy trình kiểm duyệt học thuật, thiếu tích hợp LMS và
khó gắn kết câu hỏi với chuẩn đầu ra học phần.

Luận văn đề xuất và hiện thực một hệ thống LMS mở rộng bằng Generative AI,
cho phép giảng viên đưa tài liệu vào hệ thống, tự động phân tích tài liệu,
chia nhỏ nội dung, sinh Micro-Content và Quiz, sau đó kiểm duyệt trước khi
công bố cho học viên. Hệ thống sử dụng kiến trúc Hexagonal/DDD, backend Python
gRPC, frontend Next.js, hàng đợi xử lý bất đồng bộ, PostgreSQL, MinIO, Qdrant
Vector Database và Neo4j Knowledge Graph. Lớp Curriculum-Aware phân tích đề
cương môn học, trích xuất chuẩn đầu ra học phần và liên kết chunk nội dung với
Learning Outcome để hỗ trợ sinh câu hỏi theo mục tiêu học tập.

Các thành phần AI cốt lõi gồm parser tài liệu, chunking theo heading, embedding
đa ngôn ngữ, truy xuất vector, RAG, GraphRAG định hướng thiết kế và sinh Quiz
theo Bloom taxonomy. Hệ thống được tích hợp với LMS thông qua LTI 1.3/OIDC,
trong đó Canvas LMS đóng vai trò nền tảng học tập chính, còn hệ thống AI hoạt
động như một External Tool độc lập. Vì hệ thống tuân thủ chuẩn IMS LTI,
về mặt kỹ thuật kiến trúc có thể tương thích với các LMS hỗ trợ LTI 1.3
khác, tuy nhiên việc mở rộng đó nằm ngoài phạm vi đồ án. Phần đánh giá
bao gồm kiểm thử chức năng, kiểm thử tích hợp Canvas/LTI, kế hoạch
đánh giá truy xuất, đánh giá chất lượng Quiz, hiệu năng và phân tích hạn chế.
Các số liệu chưa đo thực nghiệm được đánh dấu rõ để bổ sung trong giai đoạn
pilot.

\vspace{0.5cm}

\noindent\textbf{\large Từ khóa:} Micro-learning, Large Language Model,
Retrieval-Augmented Generation, Knowledge Graph, Bloom Taxonomy, Vector
Database, LTI, Learning Management System.
